\chapter{Analisi}

\section{Descrizione e requisiti}

Il software HearthCode ha l'obiettivo di realizzare una versione semplificata di
un gioco di carte strategico a turni ispirato a Hearthstone. L'applicazione
appartiene al dominio dei giochi in cui due partecipanti si affrontano
utilizzando risorse limitate, carte con caratteristiche specifiche e un insieme
di regole che stabilisce quali azioni possono essere eseguite in ogni momento
della partita.

A ogni partita partecipano due contendenti: il giocatore umano e un avversario
controllato automaticamente dal sistema. Entrambi dispongono di un ammontare di
punti vita, di una riserva di mana, di un mazzo da cui pescare e di una mano
contenente le carte disponibili. All'inizio di ogni turno il sistema esegue
automaticamente le operazioni obbligatorie di preparazione: aggiorna e rifornisce
il mana del giocatore attivo e pesca una carta dal suo mazzo. Se la mano ha
raggiunto la capienza massima, la carta pescata non viene aggiunta alla mano e
viene invece bruciata. In seguito il giocatore puo' schierare creature dalla
mano al campo di gioco o utilizzare quelle gia' attive per attaccare il
giocatore nemico, danneggiando le sue creature o i suoi punti vita.

Una carta puo' essere giocata solo se il giocatore possiede una quantita' di mana
almeno pari al suo costo e se il campo non ha gia' raggiunto il numero massimo
di creature consentite. Una creatura appena schierata entra in gioco in stato
inattivo e puo' attaccare solo a partire dal turno successivo del suo
proprietario. Dopo avere attaccato, la creatura viene esaurita e non puo'
compiere ulteriori attacchi nello stesso turno.

Per evitare situazioni di stallo, il gioco prevede una penalita' nel caso in cui
un giocatore debba pescare da un mazzo ormai esaurito. Poiche' i mazzi hanno la
stessa dimensione e il pescaggio e' automatico e obbligatorio per entrambi,
l'esaurimento avviene in modo bilanciato nello stesso ciclo di turni. Da quel
momento, ogni tentativo di pescaggio a mazzo vuoto danneggia i punti vita
dell'eroe con valori progressivamente crescenti, garantendo che la partita
arrivi sempre a una conclusione.

Il sistema e' responsabile della gestione dello stato della partita e
dell'applicazione coerente delle regole. Di conseguenza, deve impedire mosse non
valide, come giocare una carta non presente nella mano, usare una creatura non
ancora attiva, attaccare con una creatura gia' utilizzata o giocare una carta
senza mana sufficiente. La partita prosegue attraverso l'alternanza dei turni
fino al verificarsi di una condizione di conclusione: quando i punti vita di un
giocatore scendono a zero, l'avversario viene dichiarato vincitore.

\textbf{\\Requisiti funzionali}

\begin{itemize}
    \item Il sistema deve inizializzare una partita tra il giocatore umano e
    l'avversario controllato automaticamente, assegnando a entrambi punti vita,
    mazzo, mano iniziale e risorse di gioco.
    \item Il sistema deve gestire l'alternanza dei turni, aggiornando
    correttamente il giocatore attivo, il mana disponibile, il pescaggio
    obbligatorio e lo stato di attivazione delle creature schierate.
    \item Il sistema deve gestire il pescaggio automatico e permettere al
    giocatore di giocare carte rispettando le regole definite, in particolare il
    costo in mana, la presenza della carta nella mano e lo spazio disponibile sul
    campo.
    \item Il sistema deve bruciare la carta pescata quando la mano del giocatore
    ha raggiunto la capienza massima.
    \item Il sistema deve applicare una penalita' incrementale ai punti vita del
    giocatore quando il suo mazzo e' esaurito e deve comunque pescare una carta.
    \item Il sistema deve permettere alle creature attive di attaccare creature
    avversarie o il giocatore nemico, aggiornando punti vita, stato delle
    creature ed eventuale rimozione delle creature sconfitte.
    \item Il sistema deve impedire l'esecuzione di mosse illecite, segnalando
    l'errore senza compromettere lo stato della partita.
    \item Il sistema deve gestire autonomamente il turno dell'avversario,
    scegliendo ed eseguendo azioni valide sulla base dello stato corrente della
    partita.
    \item Il sistema deve determinare correttamente la conclusione della partita
    e il vincitore quando uno dei due giocatori esaurisce i propri punti vita.
\end{itemize}

\textbf{\\Requisiti non funzionali}

\begin{itemize}
    \item Il sistema deve essere portabile: il comportamento dell'applicazione
    deve rimanere indipendente dall'hardware e dal sistema operativo utilizzati.
    \item L'applicazione deve rispondere in tempi rapidi alle azioni dell'utente
    e ridurre i ritardi percepibili, in modo da garantire un'esperienza di gioco
    fluida.
    \item Il sistema deve mantenere uno stato di gioco consistente anche in
    presenza di input non validi, evitando che errori di interazione producano
    situazioni incoerenti.
\end{itemize}

\section{Modello del dominio}

Il sistema si basa su un tavolo di gioco, rappresentato dal \textit{board game},
che costituisce lo spazio condiviso della partita e coordina lo stato dei due
partecipanti. Il tavolo mantiene il riferimento ai giocatori, alle rispettive
armate e al turno corrente; inoltre applica le regole principali della partita,
come il posizionamento delle carte, la gestione degli attacchi, il cambio turno
e la verifica delle condizioni di vittoria.

Le carte schierate sul campo sono organizzate in due armate, una per ciascun
giocatore. L'armata rappresenta l'insieme delle creature controllate da un
partecipante e distingue le creature attive da quelle appena giocate, che devono
attendere il turno successivo prima di poter attaccare. L'armata e' inoltre
responsabile di verificare se una creatura puo' agire, di disabilitarla dopo un
attacco e di rimuoverla dal campo quando i suoi punti vita vengono azzerati.

Ogni giocatore e' caratterizzato da un identificatore, dai propri punti vita,
dalla quantita' di mana disponibile, da una mano e da un mazzo. La mano contiene
le carte che il giocatore puo' scegliere di giocare durante il turno, mentre il
mazzo rappresenta l'insieme delle carte non ancora pescate. La pesca trasferisce
una carta dal mazzo alla mano quando possibile; se il mazzo e' esaurito, il
sistema applica una penalita' ai punti vita del giocatore secondo le regole
previste.

Nella versione base del software, ogni carta giocabile e' modellata come una
creatura. Una creatura e' definita da un nome, un costo in mana, un valore di
attacco e un numero di punti vita. Il costo in mana determina se la carta puo'
essere giocata, il valore di attacco indica il danno inflitto durante un
combattimento e i punti vita determinano la permanenza della creatura sul campo.

L'avversario automatico utilizza una rappresentazione dello stato della partita
per generare le azioni disponibili, valutarne la convenienza ed eseguire una
sequenza di mosse valide durante il proprio turno. In questo modo il
comportamento dell'avversario rimane vincolato alle stesse regole applicate al
giocatore umano.
